%%%%%%%%%%%%%%%%%%%%%%%%%%%%%%%%%%%%%%%%%%%%%%%%%%%%%%%%%%%%%%
%%%%%%%%%%%%%% MA-0918 Procesos estocásticos %%%%%%%%%%%%%%%%%
%%%%%%%%%%%%%%%%%%%%%%%%%%%%%%%%%%%%%%%%%%%%%%%%%%%%%%%%%%%%%%
\newpage
\subsection*{MA-0918 Procesos estocásticos}
\addcontentsline{toc}{subsection}{MA-0918 Procesos estocásticos}

\hypertarget{MA0918}{}
\begin{refsection}
\begin{table}[H]
\centering{
\begin{tabular}{|ll|}
\hline
\textbf{Año:} IV  & \textbf{Requisitos:} MA-0840 Probabilidad \\ 
\textbf{Ciclo:} Optativa  & \textbf{Correquisitos:} Ninguno \\ 
\textbf{Tipo de curso:} Teórico & \textbf{Horas presenciales por semana:} 5 \\ 
\textbf{Créditos:} 5 & \textbf{Horas de trabajo independiente por semana:} 10 \\ \hline
\end{tabular}
}
\end{table}

\subsubsection*{Descripción del curso}

Este curso optativo introduce la teoría de procesos estocásticos en tiempo continuo, con énfasis en el movimiento browniano y las herramientas del cálculo estocástico que de él se derivan. La teoría del movimiento browniano es uno de los pilares fundamentales del estudio probabilístico moderno: desde su formulación informal en el siglo XIX, pasando por sus primeras aplicaciones en física y finanzas, hasta su formalización rigurosa a inicios del siglo XX, se ha logrado una comprensión detallada de este proceso y de sus propiedades de universalidad.

El curso se centra en el movimiento browniano, las filtraciones y martingalas asociadas, la integración estocástica, la fórmula de Itô, los procesos de Markov y las ecuaciones diferenciales estocásticas. Se supone dominio de la teoría de probabilidad basada en medida (MA-0840 Probabilidad), incluyendo convergencia de variables aleatorias, independencia y esperanza condicional.

\subsubsection*{Objetivos}

\textbf{General:} Manejar los conceptos y técnicas básicas de la teoría de procesos estocásticos en tiempo continuo, con especial atención al movimiento browniano y sus aplicaciones.

\textbf{Específicos:} Al finalizar el curso se espera que la persona estudiante sea capaz de:

\begin{enumerate}
\item Dominar los conceptos básicos de espacios y procesos gaussianos.
\item Construir el movimiento browniano y aplicar sus propiedades fundamentales.
\item Entender filtraciones, martingalas y semimartingalas continuas.
\item Utilizar el concepto de variación cuadrática en el análisis de procesos continuos.
\item Construir la integral de Itô y aplicar la fórmula de Itô en problemas relacionados.
\item Reconocer los conceptos básicos de procesos de Markov y semigrupos asociados.
\item Relacionar el movimiento browniano con ecuaciones en derivadas parciales.
\item Formular y analizar ecuaciones diferenciales estocásticas mediante resultados de existencia y unicidad.
\item Conocer nociones elementales de tiempo local del movimiento browniano.
\item Investigar y presentar oralmente un tema complementario no cubierto en el programa del curso.
\end{enumerate}

\subsubsection*{Contenidos}

\begin{enumerate}

\item \textbf{Procesos gaussianos (2 semanas):}
\begin{enumerate}
    \item Vectores y procesos gaussianos; funciones de covarianza.
    \item Ejemplos y propiedades elementales relevantes para la construcción del movimiento browniano.
\end{enumerate}

\item \textbf{Movimiento browniano (3 semanas):}
\begin{enumerate}
    \item Construcción del movimiento browniano.
    \item Propiedades básicas: continuidad, incrementos independientes y estacionarios, recurrencia y transitoriedad.
    \item Leyes del movimiento browniano; nociones introductorias de tiempo local.
\end{enumerate}

\item \textbf{Filtraciones, martingalas y semimartingalas continuas (2 semanas):}
\begin{enumerate}
    \item Filtraciones y tiempos de parada.
    \item Martingalas: teorema de parada, desigualdades clásicas.
    \item Semimartingalas continuas y descomposición canónica elemental.
\end{enumerate}

\item \textbf{Variación cuadrática e integración estocástica (2 semanas):}
\begin{enumerate}
    \item Variación cuadrática de martingalas y semimartingalas continuas.
    \item Integral de Itô: definición, propiedades y ejemplos.
\end{enumerate}

\item \textbf{Fórmula de Itô y consecuencias (2 semanas):}
\begin{enumerate}
    \item Fórmula de Itô para semimartingalas continuas.
    \item Aplicaciones: fórmula de integración por partes, teorema de Cameron--Martin--Girsanov (enunciado y usos básicos).
    \item Ecuaciones diferenciales estocásticas simples antes del tratamiento general.
\end{enumerate}

\item \textbf{Procesos de Markov (2 semanas):}
\begin{enumerate}
    \item Procesos de Markov y semigrupos.
    \item Propiedades de Markov fuerte y regularidad.
    \item Generadores infinitesimales en ejemplos relevantes.
\end{enumerate}

\item \textbf{Movimiento browniano y ecuaciones en derivadas parciales (1 semana):}
\begin{enumerate}
    \item Relación entre el movimiento browniano y la ecuación del calor.
    \item Fórmula de Feynman--Kac y aplicaciones introductorias.
\end{enumerate}

\item \textbf{Ecuaciones diferenciales estocásticas (2 semanas):}
\begin{enumerate}
    \item Existencia y unicidad de soluciones (teorema de Itô).
    \item Ejemplos y aplicaciones: modelos en finanzas, dinámica poblacional y física.
\end{enumerate}

\end{enumerate}

\subsubsection*{Metodología}

\noindent \textbf{Lineamientos metodológicos:} La persona docente a cargo de este curso debe respetar las siguientes pautas:

\begin{enumerate}
    \item El curso combina \textbf{clases magistrales} con \textbf{sesiones semanales de práctica} orientadas a la resolución de ejercicios.
    \item Las sesiones de ejercicios son participativas: la persona estudiante presenta soluciones de tareas asignadas con al menos una semana de anticipación.
    \item El trabajo independiente en casa es esencial para la comprensión de las demostraciones y la resolución de problemas.
    \item Se incluye un \textbf{proyecto individual} sobre un tema relacionado con el curso, no cubierto en el programa, con aprobación previa de la persona docente y presentación oral final.
    \item Puede utilizarse la plataforma institucional de mediación virtual para publicar material semanal y entregas; el curso admite modalidad bimodal según normativa vigente.
    \item Ante restricciones sanitarias o institucionales, las sesiones presenciales pueden trasladarse a modalidad virtual según indique la persona docente.
\end{enumerate}

\textbf{Sugerencias metodológicas:} Adicionalmente, se tienen las siguientes recomendaciones para la persona docente a cargo de este curso:

\begin{enumerate}
    \item Seguir la organización temática de \cite{LeG16} como referencia principal del curso.
    \item Alternar demostraciones formales con ejemplos que muestren la relevancia del cálculo estocástico en matemática aplicada.
    \item Promover evaluación formativa mediante tareas periódicas y exámenes parciales, además del proyecto final.
    \item Conectar los contenidos con MA-0840 Probabilidad, especialmente en convergencia, independencia y esperanza condicional.
    \item Incorporar referencias complementarias clásicas, como \cite{KS91} o \cite{RY99}, para profundización opcional.
\end{enumerate}

\nocite{LeG16}
\nocite{KS91}
\nocite{RY99}
\nocite{Oks03}

\printbibliography[heading=subbibliography,section=\the\value{refsection}]
\end{refsection}
